\documentclass[10pt]{article}
 \usepackage[margin=1in]{geometry} 
\usepackage{amsmath,amsthm,amssymb,amsfonts,color, titling}
\usepackage{listings}
\usepackage{xcolor}
\usepackage{graphicx}
\usepackage[printwatermark]{xwatermark}


\newwatermark[allpages,color=red!15,angle=45,scale=5,xpos=-12,ypos=0]{DRAFT}


\setlength{\droptitle}{-20mm} 

\usepackage[colorlinks=true,linkcolor=red,urlcolor=blue]{hyperref}

\title{Semester project instructions}
\author{PHY-905-003\\Computational Astrophysics and
  Astrostatistics\\Spring 2026}
 \date{} % leave blank to have no date

\begin{document}
 
\maketitle

\vspace{-10mm}

\section{Overall plan}

The goals of the semester project in this course are: (1) for you to
\textit{broaden and/or deepen} your understanding of the numerical
techniques you have learned in this class, (2) to apply this knowledge
to your own research or another topic of your choice in astrophysics
or physics through software \textit{that you have written} to
\textbf{analyze a dataset or model a physical situation} or
\textbf{expand on select topics we worked on in class}; and (3) to
share this knowledge with the rest of the class via a presentation and
code demonstration.  You will demonstrate achievement of these goals
through several deliverables, as detailed below.  Note that most
deadlines are given as dates, with the time of the deadline being
11:59 p.m. that day unless otherwise specified.

\textbf{In terms of project size and complexity,} this is meant to be
doable in the roughly six week time frame between now and finals week
and take roughly as much time as two homework assignments (roughly 20
total hours of work).  To that end, it may be very difficult to get
something working that directly ties to your research -- in that case,
you should choose a project that is a simplified, constrained, and/or
idealized version of what you'd actually need for your research!

Also, please note that I expect you to put all of your code, analysis
scripts, data files, and so on in the directory
\texttt{source\_code\_and\_other\_files} in this Git repository, and
regularly commit changes to the repository and push them to GitHub (in
other words, use this as the working directory for your project).
This serves three key functions: (i) it allows you to keep track of
your code changes, and revert those changes as needed; (ii) it
provides a critical backup of your project and data in case of
hardware failure, and (iii) it allows me to keep up on your progress,
and follow up as necessary.  So, please commit and push your changes
early and often!  Note also that your final code must (1) be properly
formatted and commented according to the class coding standard, and
(2) must be written in text files (not Jupyter notebooks!) and pass a
\href{https://pypi.org/project/pylint/}{PyLint} check with no
significant warnings or errors.

\vspace{3mm}
\noindent
\textbf{NOTE:} I strongly encourage you to come and talk to me at all
stages of this project -- I'm happy to brainstorm ideas, give
suggestions, and look at paper drafts and your code.  Please talk to
me before or after class or contact me via MatterMost to schedule an
appointment at another time!

\section{Project options}

There are three project options.  The first two options build directly
on topics we've learned in class, and the third is meant to be
more flexible.

\subsection{Option 1 -- Write a 2D hydro code}
\label{sec:hydro-option}

Fluid dynamics is fundamental to modeling astrophysical systems.  In
this context ``fluid dynamics'' includes both neutral fluids (as
modeled with the
\href{https://en.wikipedia.org/wiki/Euler_equations_(fluid_dynamics)}{Euler
  equations} or
\href{https://en.wikipedia.org/wiki/Navier%E2%80%93Stokes_equations}{Navier-Stokes
  equations}) and plasmas (i.e., ionized and magnetized fluids, as
modeled with the
\href{https://en.wikipedia.org/wiki/Magnetohydrodynamics}{equations of
  magnetohydrodynamics}).  It is, as a consequence, valuable to
understand how we solve these equations on a computer in order to
either use these techniques yourself or to be an educated consumer of
papers that use them.

In this project you will build on the work you've done in class
solving hyperbolic partial differential equations -- specifically,
you'll extend our second-order solution of
\href{https://en.wikipedia.org/wiki/Burgers%27_equation}{Burgers'
  equation} (the simplest nonlinear hyperbolic PDE) to solve the Euler
equations, which are the equations of inviscid hydrodynamics.  You
will do this by working through Chapters 7 and 8 of 
Michael Zingale's
\href{https://open-astrophysics-bookshelf.github.io/numerical_exercises/CompHydroTutorial.pdf}{Computational
  Astrophysical Hydrodynamics} lecture notes.  The steps you'll take
are as follows:

\begin{enumerate}

\item Read through Chapter 7 of Zingale's lecture notes and make sure
  that you understand (I) the properties of Euler's equation and (II)
  the Riemann problem as applied to these equations.

\item Implement a 1D finite volume hydro solver using the Piecewise
  Linear update (Section 8.2.2 of Zingale) and implementing the
  two-shock Riemann solver (as discussed in Section 8.3) and the
  conservative update discussed in Section 8.4.  Implement a boundary
  condition routine that allows you to set all of the boundary
  conditions described in Section 8.5, as well as periodic boundary
  conditions.  Do this assuming an ideal gas (gamma-law) equation of state.

\item Verify that your 1D hydro code works by implementing some
  standard test problems: specifically, the Sod shock tube
   and double rarefaction problem (Sec. 8.9.1), and the advection
   problem (Sec. 8.9.3).  Show that this works with the appropriate
   boundary conditions, and compare the solutions (using the L2
   norm) to either the initial conditions or reference solutions, as
   appropriate, as you vary the grid resolution and CFL condition.
   Verify also that your code works with both positive and negative
   velocities by changing the direction of the advection test and show
   that it recovers the same level of error.  Document all of this in
   plots and in text!

\item After you demonstrate that your 1D code works correctly, extend
it to 2D as described in Section 8.6 of Zingale.  Test that this code
works in each dimension individually by setting up the Sod shock tube
and double rarefaction problem in the x and (separately) y directions,
and then make sure that your code correctly couples the dimensions
using the Sedov blast wave test (Section 8.9.2).  Also use the
advection test where you allow the Gaussian profile to evolve in both
the $\pm x$ and $\pm y$ directions, as well as along all four
diagonals, and demonstrated that it works correctly and identically in
all directions.

\item Make sure to document all of your tests, including convergence
  and CFL tests, with plots and documentation.  I strongly suggest
  writing the code in a way that allows you to run it via command line
  arguments and put out plots comparing to reference solutions!

\end{enumerate}

\subsection{Option 2 -- Parallelize a 2D diffusion solver}
\label{sec:parallel-option}


\vspace{3mm}
\noindent
\textbf{\textcolor{red}{NOTE: The details of this project may change depending on
    feedback from Claire!}}
\vspace{3mm}

The two-dimensional diffusion equation is (in Cartesian coordinates
and with a constant diffusion coefficient):

\begin{equation}
\frac{\partial \phi}{\partial t} = k \left(   \frac{\partial^2 \phi}{\partial
  x^2}   +  \frac{\partial^2 \phi}{\partial y^2} \right)
\end{equation}

In this equation, $\phi$ is the quantity that is diffusing and k is
the diffusion coefficient (note that the derivation of this is
discussed in Chapter 10 of Zingale's
\href{https://open-astrophysics-bookshelf.github.io/numerical_exercises/CompHydroTutorial.pdf}{Computational
Astrophysical Hydrodynamics} lecture notes) This is a convenient
problem to experiment with because you can make it arbitrarily
computationally expensive by changing the grid resolution and/or by
changing the 3-point stencil for the second spatial derivatives to a
\href{https://en.wikipedia.org/wiki/Five-point_stencil#1D_higher-order_derivatives}{5-point
stencil}, which is a higher-order solution to the same derivative
(though going to a will require you to add an additional ghost
zone!).

The steps you'll take for this problem are:

\begin{enumerate}

  \item Implement the 2D explicit diffusion method (which is a simple
extension of the 1D version you wrote in class) entirely using numpy's
array indexing features -- in other words, don't use loops over grid
indices, as this will be important later!  Implement user-specifiable
periodic, Dirichlet, and von Neumann boundary conditions.  For initial
conditions, use the Gaussian pulse problem described in Zingale
(exercise 10.2).  Verify that you get the correct solution as a
function of time (comparing to the analytic solution using the L2
norm), and further verify that the solution in the x, y, and
annularly-averaged values are all the same.  Show that you get the
correct convergence as you increase the grid resolution as well.

\item Parallelize your 2D problem using the Python
\href{https://docs.python.org/3/library/multiprocessing.html}{multiprocessing}
module using ``tiles'' comprised of stripes of the array (i.e.,
1/N$_c$ of the domain along your y-dimension for N$_c$ CPU cores).
Using a grid whose edge (not including ghost zones) is a power of 2,
measure the simulation runtime using the Python
\href{https://docs.python.org/3/library/time.html}{time} module for
the parallel component of the problem as you increase the number of
CPU cores by powers of two.  Do this for simulation grids of various
sizes and numbers of grids, and plot the timing for both the
\href{https://hpc-wiki.info/hpc/Scaling}{strong and weak scaling}
(i.e, constant problem size vs. constant work per core), plotted
vs. ideal scaling.  You may wish to use a log-log plot for this!
(Hint: turn off all print statements for the parallel component of the
model update.)  What features do you observe in this plot, and how
might you interpret them?  (Hint: when answering this question,
consider the cost of communication and synchronization between CPU
cores, and also consider
\href{https://en.wikipedia.org/wiki/Amdahl%27s_law}{Amdahl's law}.)  I
strongly suggest running these tests on the ICER cluster, preferably
using a batch job to ensure that nobody else is running on the same
CPU cores as you.  This will ensure that your answers are more precise!

\item Parallelize your code with either \href{https://cupy.dev/}{CuPy}
  (which will work on a single GPU) or
  \href{https://mpi4py.readthedocs.io/en/stable/index.html}{mpi4py}
  (which uses message-based communication to parallelize across
  multiple CPU cores, including across multiple nodes).  CuPy has an
  excellent tutorial that you can leverage; for mpi4py, there is a
  great set of
  \href{https://github.com/jbornschein/mpi4py-examples}{mpi4py
    examples}  that you can leverage, as well as a
  \href{https://github.com/bwoshea/ACRES_MPI_tutorial}{MPI tutorial} that your
  instructor created for the ICER+CMSE REU program, which has much
  more heavily annotated versions of some of those examples.  For
  either of these examples, measure the speedup for both strong and
  weak scaling in the same way you did for a multiprocessing-based
  problem.  How do these plots differ from the one you generated in
  the previous section?  Why do you think that might be?  Note that
  you should use the ICER cluster for these experiments, and use batch
  jobs to ensure that you are not sharing the GPU or CPU resources
  with another user (which, as discussed above, will ensure that your
  answers are more precise).  In either case, also try using different
  types of hardware (V100, A100, and H200 GPUs; older-generation CPUs
  on the 2018 cluster vs. newer CPUs on the 2024 cluster) for one of
  your larger grid sizes and see if you observe any differences in
  behavior, and describe those!

  \item Make sure to document all of your tests, including convergence
  and CFL tests, with plots and documentation.  I strongly suggest
  writing the code in a way that allows you to run it via command line
  arguments and put out plots comparing to reference solutions!


\end{enumerate}

\subsection{Option 3 -- Choose your own adventure}
\label{sec:adventure}

The goal of this option is for you to apply some subset of the
numerical methods that you've learned in this course to an
astrophysical problem (or astrophysics-adjacent problem)
that's of interest to you.

One specific example of this would be to use an example problem from
your own research area (or a simplified version of a problem of that
type), determine what numerical methods are needed to solve that
problem, solve said problem numerically, and verify that the solution
behaves well numerically and makes physical sense.  For example, if
you were interested in supernova light curves you could write a simple
implicit flux-limited diffusion or Monte Carlo radiation transport
solver.  If you are interested in plasma physics, you could write a 1D
or 2D particle-in-cell (PIC) electrodynamic plasma code.  If you are
interested in understanding the light curve of an eclipsing
multi-planet solar system you could write a Bayesian fitting tool that
provides solutions (and estimates for accuracy) for fits to a light
curve.

Alternately, you could find a peer-reviewed journal article that uses
some combination of numerical methods to solve a problem that you are
interested in.  Implement your own version of the algorithm(s) that
they used to solve that problem and verify that you get the same
answers.  If you end up getting \textit{different} answers, try to
understand why and explain it.  (And appreciate that not all
peer-reviewed journal articles are correct, since the peer review
system is imperfect!)

\textbf{For your project proposal,}
briefly explain (i) the scientific motivation for the project, (ii) how
it relates to your research and/or your scientific interests, (iii) the
numerical methods that you will learn about and implement for this
project and how this goes beyond what we've learned in class in either
breadth or depth, and (iv) your intended deliverables in terms of data
analysis products or model outputs.

%\newpage
\section{Project components and deadlines}

\begin{enumerate}

\item \textbf{Project proposal, due Wednesday 3/18/2026}.  Write a
brief ($\simeq 1-2$~paragraph) proposal describing which project
option you plan to pursue.  For Options 1 and 2 just explain why
you're interested in this project.  For Option 3 (``choose your own
adventure'') please go into detail as explained in
Section~\ref{sec:adventure}.  Turn this propsal in via a text or
markdown file in the \texttt{proposal} subdirectory in this
repository.  In class the next day we'll have a brief roundtable for
everybody to explain their proposed project, and I will also follow up
with written feedback (and may ask you to modify your proposal if you
pursue Option 3 and it's overly ambitious or needs further
clarification).

\item \textbf{Project update \#1, due Monday 3/30/2026}.  The goal for
this update is for you to have done the necessary background reading
and started implementation of the code you need for your project, but
not necessarily have a fully-working code (though what you do have
should be in \texttt{source\_code\_and\_other\_files} and committed to
the repository).  The written update go in the subdirectory
\texttt{update\_1}.  For Options 1 and 2, please focus on the
challenges that you've experienced so far and the strategies you plan
to use to overcome those challenges.  For Option 3 this document
should include an update on your proposed project that details (i) the
numerical methods you're using and (if relevant) the sources you've
used to learn about them, (ii) your progress thus far, and (iii) a
discussion of any unexpected challenges you've run into and the steps
you've taken to deal with these challenges.  Note: if your project is
going poorly this is a good time to re-evaluate your project and
revise your plans, which may include changing its scope!  We will also
have a brief roundtable in class the following day for everybody to
give a quick project update, highlighting status and challenges.

\item \textbf{Project update \#2, due Wednesday 4/15/2026}.  At this
point, your code should be close to complete in terms of
implementation of its core functionality, and you should be able to
verify that it behaves as expected using simple datasets (if you've
created an analysis tool) or a simple simulation setup (if you've
written a model of some sort).  Note that these simple setups may be
useful as tests!  As with the previous update, your code should be in
the \texttt{source\_code\_and\_other\_files} directory and your
written update goes in the subdirectory \texttt{update\_2}.  For all
options the written update should describe (i) your progress thus far
and (ii) what remains to be done before your final submission.  For
Option 3, please also inclue what the tool/model is intended to do and
why you believe that it is behaving as expected based on your simpl
dataset/model setup.  \textbf{By this point, your code should also
adhere to the class coding standard and pass checks with
\href{https://pypi.org/project/pylint/}{PyLint}.}  We'll have a brief
set of updates in class the  day following this deadline.

\item \textbf{Project presentations and code demonstrations, due
before the course's final exam session the week of finals (the week of
4/27/2026 -- exact date/time for this is TBD)}.  Your final
presentation can use any of the standard presentation file formats
(Powerpoint, Keynote, Google Slides, etc.), and an exported PDF
version should be put in the \texttt{presentation\_demo} subdirectory.
Presentations will be a total of 10 minutes long, including 7-8
minutes for your talk and 2 minutes for questions and discussion.
Make sure that your presentation includes the justification for your
project and your goals (i.e., what your code is supposed to
accomplish), the numerical methods used (including a brief explanation
of how the new methods you have learned about work!), the results that
you obtained, and any lessons that you learned about the methods or
implementation that you think would be interesting to your colleagues.
If the project is an idealized, reduced, etc. version of something
you'd do for your research, please also explain that!  You should also
be prepared to give a brief demonstration of your code, if necessary.
Your presentation will be evaluated on the clarity of your visual aids
as well as the quality, clarity, and delivery of your presentation.
Please note that talk slides are due via GitHub immediately before we
start presentations.

\item \textbf{Project final code deadline, due Friday 5/1/2026}.  This
is the final deadline for your code, all of which should go in the
\texttt{final\_code} subdirectory.  \textbf{Note that this is a firm
deadline -- grades are due a few days after this, and I need time to
go through everything!}

\textbf{Your code} should: (1) be properly formatted and commented
according to the class coding standard, and (2) must pass a
\href{https://pypi.org/project/pylint/}{PyLint} check with no
\textit{significant} warnings or errors.  Comments at the top of the
main source file must describe how to run the code, including a
listing of any additional Python packages that must be installed (on
top of a baseline Miniconda distribution running a recent version of
Python 3) in order to get your code to work.  Your code will be
assessed on your implementation of the chosen algorithms, the results
that it produces, code clarity/modularity/comments, and whether it
runs on my Mac laptop and/or on ICER's clusters when using the
aforementioned Python distribution.  In addition to your code, please
include any data files that are required, as well as scripts used to
generate plots or do additional analysis.  (If the datasets are larger
than $\simeq 10$ megabytes, please talk to me before committing them
to the repository - there are other, better solutions than Git for
large files!)

\end{enumerate}

\vspace{3mm}



\end{document}
